\DocumentMetadata{
  pdfversion=2.0,
  pdfstandard=ua-2,
  lang=en-US,
  tagging=on,
  tagging-setup={math/setup=mathml-SE,tabsorder=structure}
}
\documentclass[11pt,letterpaper]{article}
\usepackage[margin=0.68in,includefoot]{geometry}
\usepackage{newcomputermodern}
\usepackage{amsmath,amscd,mathtools}
\usepackage{tikz,adjustbox}
\providecommand{\square}{\mdlgwhtsquare}
\usepackage[hidelinks]{hyperref}
\hypersetup{
  pdftitle={Complex Analysis PhD Exam, May 1992},
  pdfauthor={University of Florida Department of Mathematics},
  pdfsubject={Graduate mathematics examination},
  pdfdisplaydoctitle=true,
  bookmarksopen=true
}
\setcounter{secnumdepth}{0}
\setlength{\parindent}{0pt}
\setlength{\parskip}{0.28em}
\setlength{\emergencystretch}{3em}
\setlength{\leftmargini}{2.15em}
\setlength{\leftmarginii}{2.65em}
\setlength{\leftmarginiii}{3em}
\pagestyle{empty}
\raggedbottom
\allowdisplaybreaks
\definecolor{ProblemConcern}{HTML}{7A1F1F}
\newenvironment{problemconcern}{%
  \par\tagstructbegin{tag=Aside}\begingroup\color{ProblemConcern}
}{%
  \par\endgroup\tagstructend\nopagebreak[4]\vspace{0.12em}
}
\NewTaggingSocketPlug{block/list/label}{examproblem}{%
  \tagstructbegin{tag=Lbl}\tagstructend%
  \tagstructbegin{tag=\LBody}%
  \tagstructbegin{tag=\ProblemHeadingTag,title-o={\ProblemHeadingText}}%
  \tagmcbegin{tag=\ProblemHeadingTag,actualtext={\ProblemHeadingText}}%
  #2%
  \tagmcend%
  \tagstructend%
}
\newenvironment{examproblems}[2]{%
  \def\ProblemHeadingTag{#2}%
  \AssignTaggingSocketPlug{block/list/label}{examproblem}%
  \begin{enumerate}%
  \setcounter{enumi}{#1}%
  \renewcommand{\labelenumi}{\textbf{\theenumi.}}%
  \setlength{\topsep}{0.35em}%
  \setlength{\partopsep}{0pt}%
  \setlength{\itemsep}{0.48em}%
  \setlength{\parsep}{0.18em}%
}{\end{enumerate}}
\newenvironment{examsubparts}{%
  \AssignTaggingSocketPlug{block/list/label}{default}%
  \begin{enumerate}%
  \setlength{\topsep}{0.18em}%
  \setlength{\partopsep}{0pt}%
  \setlength{\itemsep}{0.16em}%
  \setlength{\parsep}{0.1em}%
}{\end{enumerate}}
\newenvironment{examinstructions}{%
  \par\begingroup\itshape
}{%
  \par\endgroup\vspace{0.18em}
}
\makeatletter
\renewcommand\subsection{\@startsection{subsection}{2}{0pt}%
  {-1.25ex plus -.25ex minus -.1ex}%
  {0.55ex plus .1ex}%
  {\normalfont\large\bfseries}}
\renewcommand{\@maketitle}{%
  \newpage\null\vskip -1em%
  {\footnotesize\bfseries
    \href{https://www.ufl.edu/}{University of Florida}, \href{https://math.ufl.edu/}{Department of Mathematics}\par}%
  \vskip -0.5em%
  \tagstructbegin{tag=H1,title={Complex Analysis PhD Exam, May 1992}}%
  \tagpdfparaOff%
  \tagmcbegin{tag=H1}%
    {\LARGE\bfseries\href{https://gma.math.ufl.edu/exams/complex-analysis-phd/}{Complex Analysis PhD Exam}, May 1992\par}%
  \tagmcend%
  \tagpdfparaOn%
  \tagstructend%
  \vskip 0.25em%
  \tagmcbegin{artifact}%
    \hrule height 1pt%
  \tagmcend%
  \vskip 1em%
}
\makeatother
\title{Complex Analysis PhD Exam, May 1992}
\author{}
\date{}
\begin{document}
\maketitle
\thispagestyle{empty}
\begin{examinstructions}
Do all problems. Show all the necessary details so partial credit can be properly assessed.
\end{examinstructions}
\begin{examproblems}{0}{H2}
\def\ProblemHeadingText{Problem 1.}
\item
Construct a conformal map which will take the intersection of the disks \(|z|<1\) and \(|z-1|<1\) to the unit disk.
\def\ProblemHeadingText{Problem 2.}
\item
Evaluate \(\displaystyle \int_0^\infty \frac{\log x}{1+x^4}\,dx\) using contour integration.
\def\ProblemHeadingText{Problem 3.}
\item
\begin{problemconcern}
\textbf{Warning: Suspected error.} The source reading is clear, but the hypotheses and conclusion are mathematically inconsistent: \(\sin \pi z\) has derivative~\(\pi\), not 1, at 0 and is unbounded on the imaginary axis. The intended correction is not uniquely determined, so the source text has been retained.
\end{problemconcern}
Let \(f(z)\) be entire of order 1 and the zero set of~\(f\), \(Z(f)=\{0,\pm1,\pm2,\pm3,\ldots\}\). Suppose \(f'(0)=1\), \(f\!\left(\frac12\right)=1\), and~\(f\) is bounded on the imaginary axis. Show that \(f(z)=\sin \pi z\).
\def\ProblemHeadingText{Problem 4.}
\item
Let \(P_n(z)=\displaystyle \frac{1}{n!}\frac{d^n}{dz^n}\bigl[z^n(1-z)^n\bigr]\) be the \(n^{\mathrm{th}}\)-Legendre polynomial (of degree~\(n\)). Show that \(P_n(-1)\leq (1+\sqrt{2})^{2n}\). (Hint: Write \(\displaystyle P_n(-1)=\frac{1}{2\pi i}\int_C \frac{\zeta^n(1-\zeta)^n}{(\zeta+1)^{n+1}}\,d\zeta\) over a suitable circle centered at \(-1\).)
\def\ProblemHeadingText{Problem 5.}
\item
This problem has two parts.
\begin{examsubparts}
\item[\textnormal{(a)}]
Let~\(f\) be entire and \(f(z)\longrightarrow\infty\) as \(z\longrightarrow\infty\). Show that~\(f\) is a polynomial.
\item[\textnormal{(b)}]
If in part (a) we replace~\(f\) entire by~\(f\) analytic in \(\mathbb C-\{0\}\), then what can be said?
\end{examsubparts}
\def\ProblemHeadingText{Problem 6.}
\item
Let \(f(z)\) be analytic in a domain~\(D\) and on its boundary~\(C\). If \(|f(z)|\) is constant on~\(C\), show that unless \(f(z)\) reduces to a constant there must be at least one zero of \(f(z)\) in~\(D\).
\def\ProblemHeadingText{Problem 7.}
\item
Show that the sum of the series
\[
\frac{1}{1-z}-\frac{z}{1-z^2}-\frac{z^2}{1-z^4}-\frac{z^4}{1-z^8}-\frac{z^8}{1-z^{16}}-\cdots
\]
 is 1 when \(|z|<1\), but is 0 when \(|z|>1\). This seems to contradict analytic continuation. How is this possible?
\def\ProblemHeadingText{Problem 8.}
\item
If \(f(z)\) is analytic for \(|z|<1\) and \(\displaystyle |f(z)|<\frac{1}{1-|z|}\) \((|z|<1)\), show that the coefficients \(a_n\) in the expansion
\[
f(z)=\sum_{n=0}^{\infty}a_nz^n
\]
 are subject to the inequality
\[
|a_n|\leq(n+1)\left(1+\frac1n\right)^n<e(n+1).
\]
\def\ProblemHeadingText{Problem 9.}
\item
This problem has two parts.
\begin{examsubparts}
\item[\textnormal{(a)}]
Show that the function
\[
w=\left(\frac{1+z}{1-z}\right)^{\frac{2\alpha}{\pi}},\qquad 0<\alpha<\pi,
\]
 maps \(|z|<1\) onto the angular sector \(|\arg\{w\}|<\alpha\).
\item[\textnormal{(b)}]
Prove that if \(f(z)\) is analytic in \(|z|<1\) and satisfies \(f(0)=1\) and \(|\arg\{f(z)\}|<\alpha\), then \(\displaystyle |f'(0)|<\frac{4\alpha}{\pi}\).
\end{examsubparts}
\end{examproblems}
\end{document}
